\documentclass[11pt,a4paper]{article}
\usepackage[margin=25mm]{geometry}
\usepackage{amsmath,amssymb,amsthm,booktabs,hyperref}
\newtheorem{theorem}{Theorem}[section]
\newtheorem{proposition}[theorem]{Proposition}
\newtheorem{lemma}[theorem]{Lemma}
\newtheorem{corollary}[theorem]{Corollary}
\newcommand{\gen}[1]{\langle#1\rangle}
\newcommand{\Aut}{\operatorname{Aut}}
\newcommand{\Stab}{\operatorname{Stab}}
\title{Residue stabilizers and distinguished generators}
\author{Human Compression project}
\date{14 October 2026}
\begin{document}\maketitle

This manuscript supplies the general arguments for the intake of
Fernandes--Leemans--Weiss, \emph{Highly symmetric hypertopes} (2016).
Its central observation is that a stabilizer of a flag can be read from
paths of chambers in its residue.  Ordinary paths give the distinguished
involutions of a regular hypertope; even paths give the distinguished
generators of a chiral hypertope.  Intersections then follow from
intersections of pointwise stabilizers.  We use the already acquired
propagation, extension-fiber and reflection theorems in the chamber
manuscript.  Their proofs are dependencies, not new assertions of the
present application.  Finite examples, numerical tables and the toroidal
families are separate components.

\section{Conventions and acquired geometric arguments}

Let $I$ be a finite type set.  An incidence geometry has one object of each
type in every chamber, and every flag extends to a chamber.  Distinct
objects of the same type are not incident.  The residue of a flag $F$
consists of the objects of the remaining types incident with every object
of $F$.  Residual connectedness means incidence connectedness of every
residue of rank at least two, including the empty-flag residue.  Thinness
means that every panel is contained in exactly two chambers.

For a thin geometry write $C^i$ for the unique chamber different from $C$
on the panel obtained by deleting its type-$i$ object.  Superscripts are
applied in their written order: $C^{ij}=(C^i)^j$.  All group actions in this
manuscript are \emph{left} actions; thus $(gh)C=g(hC)$.  An automorphism
preserves types unless stated otherwise.  It commutes with every operation
$C\mapsto C^i$.

The acquired propagation and reconstruction lemma has three conclusions:
\begin{enumerate}
\item Chambers containing a fixed flag are connected by adjacency steps
which retain that flag.
\item In a thin residually connected geometry, an automorphism fixing a
chamber is the identity.
\item A permutation of chambers preserving every adjacency color comes
from a unique automorphism of the incidence geometry.
\end{enumerate}
The proof of the first conclusion is an induction on the rank of the
residue.  It is important to apply the induction to residues known to be
connected.  Mere incidence connectedness of the whole geometry does not
supply that hypothesis.  The converse needed for Proposition~2.1 of the
source is elementary: a chamber path in a residue of rank at least two
gives an incidence path between any two of its objects.  Consecutive
chambers share at least one object, and each object lies in a chamber.
Thus residual connectedness and strong chamber connectedness are
equivalent under the source's hypotheses.  This corrects an intermediate
step in its proof without changing the stated proposition.

Chamber transitivity implies transitivity on flags of every fixed type:
extend both flags to chambers and carry one chamber to the other.
Conversely chamber transitivity is a special case of flag transitivity.
This is Proposition~2.2 of the source.

\section{The residue stabilizer identity}

Fix a chamber $C=\{c_i:i\in I\}$ and put
\[
             F_{I\setminus S}=\{c_i:i\notin S\}
             \qquad(S\subseteq I).
\]
Regularity means transitivity of the full automorphism group on chambers;
by propagation its action is free as well.  The same conclusions below
hold for any chamber-transitive group of actual automorphisms, which then
equals the full group by freeness.

\begin{theorem}\label{thm:regular}
Let $G$ act regularly on the chambers of a hypertope.  There is a unique
$\rho_i\in G$ with $\rho_iC=C^i$ for each $i\in I$.  These elements are
distinct nonidentity involutions.  For every $S\subseteq I$,
\begin{equation}\label{eq:regular-stabilizer}
       \gen{\rho_i:i\in S}=\Stab_G(F_{I\setminus S}).
\end{equation}
In particular they generate $G$ and satisfy the full intersection property
\[
 \gen{\rho_i:i\in S}\cap\gen{\rho_i:i\in T}
                  =\gen{\rho_i:i\in S\cap T}.
\]
\end{theorem}
\begin{proof}
Existence and uniqueness of $\rho_i$ follow from regularity.  Since
$\rho_i(C^i)=(\rho_iC)^i=C$, its square fixes $C$, so $\rho_i^2=1$.
It is nonidentity because $C^i\ne C$.  The chambers $C^i,C^j$ are distinct
when $i\ne j$, so the generators are distinct.  Each $\rho_i$ fixes
every $c_j$ with $j\ne i$.  This proves inclusion from left to right in
\eqref{eq:regular-stabilizer}.

Let $g$ fix $F_{I\setminus S}$.  The chambers $C$ and $gC$ contain this
flag.  Propagation inside its residue gives a path from $C$ to $gC$ whose
colors all belong to $S$.  If its colors are $i_1,\ldots,i_m$, its
successive chambers are
\[
 C,\quad \rho_{i_1}C,\quad
 \rho_{i_1}\rho_{i_2}C,\quad\ldots,\quad
 \rho_{i_1}\cdots\rho_{i_m}C.
\]
Indeed adjacency commutes with every group element.  The last product and
$g$ have the same image of $C$, and freeness makes them equal.  This also
works for $S=\varnothing$, when the stabilizer is trivial.  It proves
\eqref{eq:regular-stabilizer}.  Taking $S=I$ proves generation.
Finally, to fix both complementary flags is exactly to fix their union,
which is the complementary flag of $S\cap T$.  Applying the identity
three times gives the claimed intersection.
\end{proof}

This proves source Theorem~4.1 without an induction that presumes the
identification of the residue group.  Its rank-two Coxeter quotient
interpretation follows as well: put $m_{ij}=o(\rho_i\rho_j)$, allowing
$\infty$.  The Coxeter presentation with these labels maps onto $G$.
The rank-two subgroups are dihedral of order $2m_{ij}$ (infinite when the
label is infinite), so the map is injective on each such subgroup.

For the chiral case, $G=\Aut_I(\Gamma)$ has exactly two chamber orbits,
and adjacent chambers lie in different orbits.  Call the orbit of $C$
positive.  Every two-step chamber $C^{ij}$ is positive.  The action is
free by propagation, so there is a unique element taking $C$ to $C^{ij}$.
Fix $0\in I$, put $\alpha_0=1$, and for $i\ne0$ define
$\alpha_iC=C^{0i}$.

\begin{theorem}\label{thm:chiral}
For a chiral hypertope of rank at least three set
\[
 \alpha_{ij}=\alpha_i^{-1}\alpha_j,\qquad
 H_S=\gen{\alpha_{ij}:i,j\in S}\quad(S\subseteq I).
\]
Then $\alpha_{ij}C=C^{ij}$ and
\begin{equation}\label{eq:chiral-stabilizer}
                 H_S=\Stab_G(F_{I\setminus S}).
\end{equation}
Consequently $G=\gen{\alpha_i:i\ne0}$ and
$H_S\cap H_T=H_{S\cap T}$ for all subsets $S,T$.  The set of nonidentity
distinguished generators $\{\alpha_i:i\ne0\}$ is independent.
\end{theorem}
\begin{proof}
First $\alpha_i^{-1}C=C^{i0}$, since applying $\alpha_i$ to the latter
chamber gives $C^{0ii0}=C$.  Therefore
\[
 \alpha_i^{-1}\alpha_jC
       =\alpha_i^{-1}(C^{0j})=(C^{i0})^{0j}=C^{ij}.
\]
The same identity includes $i=0$, $j=0$ and $i=j$.  In particular
$\alpha_{ij}$ fixes every object of $C$ outside types $i,j$, so
$H_S$ is contained in the stabilizer in \eqref{eq:chiral-stabilizer}.

For the reverse inclusion let $g$ fix the complementary flag.  Propagation
gives a path from $C$ to $gC$ using colors in $S$.  Every adjacency changes
the chamber orbit.  Since both endpoints are positive, the path has even
length, say with colors $i_1,j_1,\ldots,i_m,j_m$.  After its first pair
the chamber is $\alpha_{i_1j_1}C$.  Equivariance then shows successively
that its last chamber is
$\alpha_{i_1j_1}\cdots\alpha_{i_mj_m}C$.
Freeness makes this product equal to $g$.

This argument includes singleton and empty $S$.  Their group $H_S$ is
trivial; equivalently a group element stabilizing a panel cannot exchange
its two chambers because those chambers have different colors.  We have
proved the stabilizer identity for every $S$, rather than assumed that a
residue has at most two orbits.  Generation follows on taking $S=I$;
intersections follow by taking the union of the two fixed flags.

The chamber $C^{0i}$ differs from $C$ for $i\ne0$: equality would imply
$C^0=C^i$, which is impossible for different types.  Hence $\alpha_i\ne1$.
The subgroup generated by the other $\alpha_j$ is
$H_{I\setminus\{i\}}$, whereas $\alpha_i\in H_{\{0,i\}}$.  Their
intersection is $H_{\{0\}}=1$.  Thus no $\alpha_i$ lies in the subgroup
generated by the others.
\end{proof}

The intersection conclusion contains the source's condition $IC^+$, which
is stated only for $|S|,|T|\ge2$.  This proves Theorem~7.1 and
Corollary~7.2.  Ranks at most two have no chiral hypertopes: a connected
thin rank-two incidence graph is a cycle or a doubly infinite path, whose
type-preserving automorphisms are transitive on its edges; the rank-one
case is the two-object geometry.

\begin{corollary}\label{cor:properflags}
The chiral group is transitive on flags of every proper type set.
In particular every proper truncation is flag-transitive.
\end{corollary}
\begin{proof}
Let $F$ be a flag of proper type $J$ and extend it to a chamber $D$.
If $D$ is negative, choose $i\notin J$ and replace it by $D^i$; this
retains $F$ and makes the chamber positive.  Some $g\in G$ takes $C$
to the resulting positive chamber and therefore takes its $J$-subflag
to $F$.  A flag in a truncation extends to a chamber of that truncation
by extending it in the whole geometry.  The same transitivity argument
applies there.
\end{proof}
This includes Proposition~8.1 of the source.  It asserts transitivity,
not thinness or regularity of those truncations.

\section{Cosets and the converse construction}

For subgroups $P_i\le G$, the objects of type $i$ are left cosets
$gP_i$, with incidence given by nonempty intersection.  This always
defines an incidence system and has the base chamber $\{P_i:i\in I\}$.
It need not have the flag-extension property in rank at least four.
Left multiplication is transitive on objects of each type and on
incident pairs: take a common representative of the two cosets.
These are the precise assertions used from source Proposition~3.1.

The kernel of the action is
\[
              K=\bigcap_{i\in I}\bigcap_{g\in G}gP_ig^{-1}.
\]
It is contained in every $P_i$.  Sending $gP_i$ to
$(gK)(P_i/K)$ is a bijection on each type.  It preserves incidence in
both directions, because a representative of an intersecting quotient
coset can be lifted and the discrepancy is in $K\le P_i$.
Thus it gives the isomorphism of source Theorem~3.3.  A claim that $G$
itself acts faithfully before this quotient would be false.

If a group is transitive on objects and incident pairs of each type,
its geometry is reconstructed by these cosets of the stabilizers of a
base chamber.  Indeed $gP_i\mapsto gc_i$ is a well-defined bijection;
two image objects are incident precisely when the pair is a translate
of the base pair, which is equivalent to the two cosets intersecting.
This applies to regular hypertopes and, by Corollary~\ref{cor:properflags},
to chiral hypertopes of rank at least three.  The residue stabilizer
identities identify their $P_i$ with the prescribed marked subgroups.

For a regular faithful action the Borel subgroup
$B=\bigcap_iP_i$ is the chamber stabilizer and hence is trivial.  The
stabilizer of a panel has exactly two elements, since it acts transitively
and freely on the two containing chambers.  These observations prove
Lemma~3.4 and its stated consequence.  For a nonfaithful chamber-transitive
coset action on a hypertope, the corresponding conclusion is $B=K$ and
a panel stabilizer has quotient of order two by $K$; passage to the
faithful quotient is essential.

\begin{theorem}\label{thm:converse}
Let $(G,\{\rho_i:i\in I\})$ be a C-group and put
$G_S=\gen{\rho_i:i\in S}$ and $P_i=G_{I\setminus\{i\}}$.
If its coset incidence system is flag-transitive under $G$, it is a
regular hypertope.  Its full automorphism group is $G$.
\end{theorem}
\begin{proof}
Every flag is a translate of the corresponding base flag, so extends
to a chamber.  Iterated use of the C-group intersection property gives
\[
                  \bigcap_{j\notin S}P_j=G_S
                  \qquad(S\subseteq I),
\]
including the conventions for empty intersections.  The base chamber
stabilizer is $G_\varnothing=1$, so the action is faithful and free
on chambers.  The stabilizer of the complementary base flag acts
transitively on the chambers containing that flag: a group element taking
the base chamber to any such chamber must fix every object of the flag.
For a panel its stabilizer is $G_{\{i\}}=\{1,\rho_i\}$; hence the
geometry is thin.

More generally its residue chambers are the chambers $gC$ with
$g\in G_S$.  A word in the $\rho_i$, $i\in S$, gives an adjacency path
between $C$ and $gC$ wholly inside that residue.  Thus all residues
are chamber-connected and therefore incidence-connected in rank at
least two.  The geometry is a hypertope.  Its full automorphism group
acts freely on chambers, whereas $G$ already acts transitively, so the
two groups coincide.
\end{proof}
This proves source Theorem~4.6.  It does not assert flag transitivity
from the intersection property alone.

\section{The classical coset criteria}

The following elementary arguments record exactly which group conditions
are used.  Their origins remain Tits, Dehon and Buekenhout--Hermand as
attributed in the source; they are not new claims of priority.

\begin{proposition}\label{prop:dehon}
Suppose the coset incidence system is flag-transitive.  Put
$P_J=\bigcap_{j\in J}P_j$, with $P_\varnothing=G$.
It is residually connected if and only if
\[
       P_J=\gen{P_{J\cup\{k\}}:k\in I\setminus J}
                    \qquad(|I\setminus J|\ge2).
\]
\end{proposition}
\begin{proof}
Fix the base flag of type $J$ and write $H=P_J$.
Every incident object of type $k\notin J$ has the form $hP_k$ with
$h\in H$, by flag transitivity on $J\cup\{k\}$.  Two such representatives
give the same object precisely when they give the same left coset of
$H\cap P_k$ in $H$.  If two objects of the residue are incident, flag
transitivity on their union with the base flag supplies a common
representative in $H$.  Consequently the residue is exactly the coset
incidence system of $H$ with subgroups $H\cap P_k$.

In a coset incidence graph of rank at least two, the component containing
the base flag is invariant under each of its vertex stabilizers: an
automorphism fixing a vertex preserves that vertex's component.  It
therefore contains every vertex represented by an element of the group
$L$ generated by those stabilizers.  Conversely, since each stabilizer
is contained in $L$, assigning the label $hL$ to a vertex represented by
$h$ is well-defined and constant along edges.  The graph is connected
if and only if $L=H$.  Apply this to every $J$.
\end{proof}
Without flag transitivity, a residue can have additional incidence edges;
the equality with a coset graph must not be assumed.

\begin{proposition}\label{prop:flagtest}
For each nonempty $J\subseteq I$ choose $a(J)\in J$.
The coset system is flag-transitive if and only if, for every $|J|\ge3$,
\[
 \bigcap_{j\in J\setminus\{a(J)\}}P_jP_{a(J)}
       =\left(\bigcap_{j\in J\setminus\{a(J)\}}P_j\right)P_{a(J)}.
\]
\end{proposition}
\begin{proof}
The acquired extension-fiber theorem identifies orbits on flags of type
$T\cup\{k\}$, assuming transitivity on type $T$, with
\[
 P_T\backslash\left(\bigcap_{j\in T}P_jP_k\right)/P_k.
\]
There is one orbit precisely when the middle set is $P_TP_k$.
Transitivity on flags of size at most two always holds.  Inducting on
$|J|$, with $T=J\setminus\{a(J)\}$, proves sufficiency.  Necessity
is the same orbit calculation at every step.  Every flag then translates
to a base flag and extends, so no missing geometry hypothesis is hidden
in this test.
\end{proof}

\begin{lemma}\label{lem:helly}
For subgroups $A,B,D$ the following are equivalent:
\begin{align*}
 AB\cap AD&=A(B\cap D),\tag{i}\\
 (A\cap B)(A\cap D)&=A\cap BD,\tag{ii}
\end{align*}
and (iii): three right cosets of $A,B,D$ that meet pairwise have a
common point.
\end{lemma}
\begin{proof}
Suppose (i), and let $x=bd\in A\cap BD$.  Then $b=xd^{-1}$ belongs
to $AB\cap AD$, so $b=ae$ with $a\in A$ and $e\in B\cap D$.
Since $b,e\in B$, we have $a\in A\cap B$; and $ed=a^{-1}x$
belongs to $A\cap D$.  This proves the nontrivial inclusion of (ii).

Conversely let $x=ab=cd\in AB\cap AD$, with $a,c\in A$,
$b\in B$, $d\in D$.  Inverting equality (ii) also gives
$A\cap DB=(A\cap D)(A\cap B)$.  Thus
$c^{-1}a=db^{-1}=uv$, where $u\in A\cap D$ and $v\in A\cap B$.
It follows that $x=cu(vb)$, with $vb=u^{-1}d\in B\cap D$.
This gives (i); the opposite inclusions in both equalities are automatic.

For the coset assertion, translate a common point of the $B$- and
$D$-cosets to the identity by right multiplication.  The three cosets
then have the form $Ak,B,D$.  They meet pairwise precisely when
$k\in AB\cap AD$.  They have a common point precisely when
$k\in A(B\cap D)$.  This is exactly (i).
\end{proof}
This is source Lemma~4.2.  Inverting cosets gives the equivalent statement
for the left-coset convention used elsewhere here.

\begin{corollary}\label{cor:rankthree}
For a rank-three C-group the standard coset geometry is thin if and only
if the designated group acts regularly.  In either case it is a regular
hypertope.
\end{corollary}
\begin{proof}
Any rank-three coset system is a geometry: an incident pair has a common
representative, which supplies an object of the third type.  If the group
is chamber-transitive, Theorem~\ref{thm:converse} applies.

Now suppose it is thin.  Normalize any incident pair to $P_i,P_j$,
and let $k$ be the remaining type.  There are two distinct completions
$P_k$ and $\rho_kP_k$: the element $\rho_k$ belongs to $P_i\cap P_j$
and not to $P_k$, by the intersection property.  Both completed chambers
are translates of the base chamber.  Thinness allows no other completion.
Hence every chamber is a translate of the base chamber, and the theorem
applies.  No inference of freeness on all chambers from the base stabilizer
alone is used.
\end{proof}

\section{Rotation subgroups and strings}

\begin{proposition}\label{prop:rotation}
Let $(G,\{\rho_i:i\in I\})$ be a C-group with $0\in I$.
Set $\alpha_0=1$, $\alpha_i=\rho_0\rho_i$ for $i\ne0$ and
$H=\gen{\alpha_i:i\ne0}$.  These generators are independent.
The index $[G:H]$ is one or two.  If it is two, then
\[
 \gen{\alpha_i^{-1}\alpha_j:i,j\in S}=G_S\cap H
                         \qquad(S\subseteq I),
\]
so the rotation generators satisfy $IC^+$.
\end{proposition}
\begin{proof}
Conjugation by $\rho_0$ inverts every $\alpha_i$, and
$\rho_i=\rho_0\alpha_i$.  Thus $H\lhd G$ and $G/H$ is generated
by the image of $\rho_0$, whose order is at most two.
If $\alpha_i$ lay in the subgroup generated by the other $\alpha_j$,
it would belong to $G_{I\setminus\{i\}}$.  Multiplication by
$\rho_0$ would put $\rho_i$ in that subgroup, contradicting the
intersection property.  This proves independence regardless of the index.

If the index is two, all $\rho_i$ have the same nontrivial image in
$G/H$.  Consequently a word in the $\rho_i$, $i\in S$, represents
an element of $H$ exactly when its length is even.  Every such word is a
product of pairs $\rho_i\rho_j=\alpha_i^{-1}\alpha_j$ with
$i,j\in S$.  This proves the displayed identity.  Intersecting two
such identities and using $G_S\cap G_T=G_{S\cap T}$ proves $IC^+$.
\end{proof}
These are Propositions~6.1 and~6.2.  The index-two hypothesis in the last
assertion cannot be discarded; the source's Example~6.3 is checked in the
finite component.

For source Theorem~5.1 we retain the standard string C-group correspondence
with abstract regular polytopes, in its exact marked form: the faces of
type $i$ are the cosets of $G_{I\setminus\{i\}}$, its flag action is
regular, and its rank-two sections have polygon sizes
$o(\rho_i\rho_j)$.  This is the classical Aschbacher--Schulte input
cited in that theorem and the McMullen--Schulte correspondence already
retained by the programme.  It yields flag transitivity, after which
Theorem~\ref{thm:converse} gives the hypertope statement.  This input is
not counted as a new proof here.

Conversely, for a regular hypertope the rank-two residue on types $i,j$
has group $G_{\{i,j\}}$ by Theorem~\ref{thm:regular}.  Its alternating
chamber cycle has length $2o(\rho_i\rho_j)$, or is infinite if that
order is infinite.  A missing edge in its diagram means this residue is
a digon, so the product order is two.  If the diagram is a string, every
pair of nonconsecutive types is missing, and the distinguished C-group
is therefore a string C-group.  This proves source Theorem~5.2.

\section{The two chamber orbits in Construction 8.1}

Let $I=\{0,\ldots,r-1\}$ with $r\ge3$.  Suppose
$G=\gen{\alpha_1,\ldots,\alpha_{r-1}}$ is a $C^+$-group with its
independent nonidentity generators.  Set $\alpha_0=1$ and
$H_S=\gen{\alpha_i^{-1}\alpha_j:i,j\in S}$.
The parabolics of Construction~8.1 are
\[
 P_i=H_{I\setminus\{i\}}\quad(i\in I).
\]
For $i>0$ this is $\gen{\alpha_j:j\ne0,i}$; for $i=0$ it is
$\gen{\alpha_1^{-1}\alpha_j:2\le j<r}$, by telescoping the
pairwise differences.  We continue to assume, as the source theorem does,
that the associated coset system is a hypertope and that $G$ has
\emph{exactly two} orbits on its chambers.

Iterating $IC^+$ gives trivial intersections of all parabolics but one:
the common support is a singleton and its $H$-group is trivial.
The iteration uses $IC^+$ while both supports have size at least two;
as soon as a singleton is reached the intersection is already trivial.
The base chamber stabilizer is trivial as well.  Thus this is a faithful
action, and propagation makes it free on every chamber.

Write $C=\{P_i:i\in I\}$.  The coset $\alpha_1P_0$ contains every
$\alpha_j$, $j>0$.  For each $i>0$ choose $j>0$ different from $i$;
then $\alpha_j\in P_i$, so replacing $P_0$ by $\alpha_1P_0$ gives
a chamber.  It differs from $C$: otherwise
$\alpha_1\in H_{I\setminus\{0\}}\cap H_{\{0,1\}}=H_{\{1\}}=1$.
By thinness it is $C^0$; denote it by $C_0$.

For $i>0$, the coset $\alpha_i^{-1}P_i$ is different from $P_i$
by independence.  It meets $P_k$ for $k>0$, $k\ne i$, at
$\alpha_i^{-1}$.  It meets $P_0$ at $\alpha_i^{-1}\alpha_j$,
where $j>0$, $j\ne i$.  Replacing $P_i$ by this coset therefore
gives the chamber $C^i$.  On multiplying it by $\alpha_i$, every
positive-type object becomes the corresponding object of $C_0$, while
$\alpha_iP_0=\alpha_1P_0$.  Hence
\[
                         C^i=\alpha_i^{-1}C_0
                         \qquad(i\in I).
\]
Here $i=0$ follows from $\alpha_0=1$.

No element of $G$ takes $C$ to any $C^i$, since it would fix the base
panel and its stabilizer is trivial.  The same holds at any translate
of $C$.  Also every neighbor of a translate of $C_0$ is a translate
of $C$: this follows from $C^i=\alpha_i^{-1}C_0$ and involutivity
of adjacency.  With the exactly-two-orbits hypothesis, all chambers
are therefore uniquely
\[
       (g,+)=gC,\qquad(g,-)=gC_0,
       \qquad (g,+)\sim_i(g\alpha_i^{-1},-).
\]
These are precisely the hypotheses and coordinates of the acquired
two-regular-orbits theorem.  It follows that the full group is $G$, and
the hypertope is chiral, unless some $\theta\in\Aut(G)$ inverts every
$\alpha_i$.  If such a $\theta$ exists, the explicit color-preserving map
\[
       (g,+)\mapsto(\theta(g),-),\qquad
       (g,-)\mapsto(\theta(g),+)
\]
extends to an incidence involution $F$, and the full group is
$G\rtimes\gen{F}$, with $FgF^{-1}=\theta(g)$.  This proves source
Theorem~8.2 by the already acquired theorem.  In this notation the
adjoined chamber involution can be nontrivial even when the abstract
automorphism $\theta$ is the identity.

The proof uses the two-orbits hypothesis to identify the entire chamber
set.  It does not establish that hypothesis for an arbitrary $C^+$
presentation.  Nor does an abstract semidirect product alone establish
that its extra involution acts on the proposed geometry.

\section{Scope, attribution and source questions}

Theorems~\ref{thm:regular} and~\ref{thm:chiral} replace the two separate
intersection arguments by the same residue-stabilizer identity.
The converse construction and the rotation argument are then short
deductions.  The geometric propagation lemma, extension-fiber theorem
and explicit reflection theorem are acquired project proofs, with their
earlier attributions unchanged.  The subgroup criteria remain the
classical results attributed by the source to Tits, Dehon and
Buekenhout--Hermand.  The string correspondence remains an explicit
retained theorem, not an inferred consequence of an order calculation.
No claim of historical priority for the stabilizer observation is made.

The programme already has reviewed negative answers to source
Problems~9.3 and~9.4 under the standard parabolic construction.  Their
precise statements are retained here.  Thinness and chamber connectedness
force the standard C-group orbit to contain all chambers, excluding the
example requested in~9.3.  Every rank-three standard coset system on
three generators is residually connected, excluding the example requested
in~9.4.  The latter argument only needs a connected \emph{spanning}
coset graph in each residue and does not assert an equality of residue
incidence relations without flag transitivity.  These are acquisitions
of earlier answers, not two new solutions generated by this intake.

Problems~9.1 and~9.2 and the conjectured completeness of the known chiral
list remain separately registered questions.  This component supplies no
classification of all locally toroidal cases.  It grants no whole-original
credit before the finite and family components and the final source map
are reviewed.  Existing Lean declarations retain exactly their previous
algebraic and action scope; no new formalization or compression measurement
is asserted here.

\begin{thebibliography}{9}
\bibitem{FLW} M. E. Fernandes, D. Leemans and A. Ivi\'c Weiss,
\emph{Highly symmetric hypertopes}, Aequationes Mathematicae 90 (2016),
1045--1067. Fixed accepted author manuscript, 22 pages.
\bibitem{HC} Human Compression project, \emph{Chambers and marked covers},
topic-chamber/v2, sections on normalization, propagation and reconstruction,
and two regular orbits. The accepted version and review are acquired by
their exact artifact bindings.
\bibitem{MS} P. McMullen and E. Schulte, \emph{Abstract Regular Polytopes},
Cambridge University Press, 2002, Section 2E. The marked regular-polytope
correspondence is retained, as in the existing programme.
\end{thebibliography}
\end{document}
